\documentclass[10pt,a4paper]{amsart}
\usepackage{amsmath,amssymb,amsthm,graphicx}
\usepackage{verbatim}
\newenvironment{reference}{\comment}{\endcomment}
%\newenvironment{reference}{}{}
\newenvironment{slogan}{\comment}{\endcomment}
\newenvironment{history}{\comment}{\endcomment}

% For commutative diagrams we use Xy-pic
\usepackage[all]{xy}

% We use 2cell for 2-commutative diagrams.
\xyoption{2cell}
% Theorem environments.
%
\theoremstyle{plain}
\newtheorem{theorem}[subsection]{Theorem}
\newtheorem{proposition}[subsection]{Proposition}
\newtheorem{lemma}[subsection]{Lemma}

\theoremstyle{definition}
\newtheorem{definition}[subsection]{Definition}
\newtheorem{example}[subsection]{Example}
\newtheorem{exercise}[subsection]{Exercise}
\newtheorem{situation}[subsection]{Situation}

\theoremstyle{remark}
\newtheorem{remark}[subsection]{Remark}
\newtheorem{remarks}[subsection]{Remarks}

\numberwithin{equation}{subsection}

% Macros
%
\def\lim{\mathop{\mathrm{lim}}\nolimits}
\def\colim{\mathop{\mathrm{colim}}\nolimits}
\def\Spec{\mathop{\mathrm{Spec}}}
\def\Hom{\mathop{\mathrm{Hom}}\nolimits}
\def\Ext{\mathop{\mathrm{Ext}}\nolimits}
\def\SheafHom{\mathop{\mathcal{H}\!\mathit{om}}\nolimits}
\def\SheafExt{\mathop{\mathcal{E}\!\mathit{xt}}\nolimits}
\def\Sch{\mathit{Sch}}
\def\Mor{\mathop{\mathrm{Mor}}\nolimits}
\def\Ob{\mathop{\mathrm{Ob}}\nolimits}
\def\Sh{\mathop{\mathit{Sh}}\nolimits}
\def\NL{\mathop{N\!L}\nolimits}
\def\CH{\mathop{\mathrm{CH}}\nolimits}
\def\proetale{{pro\text{-}\acute{e}tale}}
\def\etale{{\acute{e}tale}}
\def\QCoh{\mathit{QCoh}}
\def\Ker{\mathop{\mathrm{Ker}}}
\def\Im{\mathop{\mathrm{Im}}}
\def\Coker{\mathop{\mathrm{Coker}}}
\def\Coim{\mathop{\mathrm{Coim}}}

% Boxtimes
%
\DeclareMathSymbol{\boxtimes}{\mathbin}{AMSa}{"02}

%
% Macros for moduli stacks/spaces
%
\def\QCohstack{\mathcal{QC}\!\mathit{oh}}
\def\Cohstack{\mathcal{C}\!\mathit{oh}}
\def\Spacesstack{\mathcal{S}\!\mathit{paces}}
\def\Quotfunctor{\mathrm{Quot}}
\def\Hilbfunctor{\mathrm{Hilb}}
\def\Curvesstack{\mathcal{C}\!\mathit{urves}}
\def\Polarizedstack{\mathcal{P}\!\mathit{olarized}}
\def\Complexesstack{\mathcal{C}\!\mathit{omplexes}}
% \Pic is the operator that assigns to X its picard group, usage \Pic(X)
% \Picardstack_{X/B} denotes the Picard stack of X over B
% \Picardfunctor_{X/B} denotes the Picard functor of X over B
\def\Pic{\mathop{\mathrm{Pic}}\nolimits}
\def\Picardstack{\mathcal{P}\!\mathit{ic}}
\def\Picardfunctor{\mathrm{Pic}}
\def\Deformationcategory{\mathcal{D}\!\mathit{ef}}
\usepackage{cleveref}
\crefname{originallemma}{original lemma}{original lemmas}
\Crefname{originallemma}{Original lemma}{Original lemmas}
\makeatletter
\newlabel{hypercovering-lemma-hypercovering-continuous-functor}{{0DAZ}{source}{}{sourcepointer-1}{}}
\newlabel{hypercovering-lemma-hypercovering-continuous-functor@cref}{{[originallemma][0][]0DAZ}{[1][0][]source}{}{}{}}
\makeatother
\crefname{originallemma}{original lemma}{original lemmas}
\Crefname{originallemma}{Original lemma}{Original lemmas}
\makeatletter
\newlabel{hypercovering-lemma-w-contractible}{{094K}{source}{}{sourcepointer-2}{}}
\newlabel{hypercovering-lemma-w-contractible@cref}{{[originallemma][0][]094K}{[1][0][]source}{}{}{}}
\makeatother
\crefname{originalremark}{original remark}{original remarks}
\Crefname{originalremark}{Original remark}{Original remarks}
\makeatletter
\newlabel{hypercovering-remark-take-unions-hypercovering-X}{{0DB2}{source}{}{sourcepointer-3}{}}
\newlabel{hypercovering-remark-take-unions-hypercovering-X@cref}{{[originalremark][0][]0DB2}{[1][0][]source}{}{}{}}
\makeatother
\crefname{originalreference}{original reference}{original references}
\Crefname{originalreference}{Original reference}{Original references}
\makeatletter
\newlabel{lemma-get-many-weakly-contractible}{{098I}{source}{}{sourcepointer-4}{}}
\newlabel{lemma-get-many-weakly-contractible@cref}{{[originalreference][0][]098I}{[1][0][]source}{}{}{}}
\makeatother
\crefname{originalreference}{original reference}{original references}
\Crefname{originalreference}{Original reference}{Original references}
\makeatletter
\newlabel{lemma-weakly-contractible-cover}{{0F4P}{source}{}{sourcepointer-5}{}}
\newlabel{lemma-weakly-contractible-cover@cref}{{[originalreference][0][]0F4P}{[1][0][]source}{}{}{}}
\makeatother
\title{Every object of the pro-etale site has a hypercovering whose terms are single weakly contractible objects}
\author{The Stacks Project Authors}
\date{stacks-project:proetale.tex:local-snapshot-20261012; GNU FDL 1.2 or later}
\begin{document}
\maketitle
\noindent\textit{Editorial note.} All delivered mathematical text is the author's own, copied from the pinned Stacks Project chapter file proetale.tex; no statement or proof step was rewritten, reordered or removed.\par\medskip
\noindent\textit{Editorial note.} The counted proof is the author's complete argument: the application of the hypercovering construction to finite coverings by weakly contractible objects, the passage to degreewise disjoint unions, and the repetition of the argument inside the subcategory of quasi-compact and quasi-separated objects with affine weakly contractible coverings.\par\medskip
\noindent\textit{Editorial note.} The label lemma-w-contractible-hypercovering names the counted statement, and the proof environment immediately following it in the source is delivered verbatim as its proof.\par\medskip
\noindent\textit{Editorial note.} The counted statement is carried by the author's own lemma environment in the chapter source.\par\medskip
\noindent\textit{Editorial note.} Pre-open gate: python operations/agent-authoring/tools/precard\_gate.py --source proetale.tex --label lemma-w-contractible-hypercovering --start 3215 --end 3266 --local-id 3579 returned RESULT=OK (L3579 free, occupancy FREE, public interval clear). Candidate taken from sources/stacks-proetale/candidates-halfA.json.\par\medskip
\noindent\textit{Editorial note.} The statement and proof are copied from the pinned chapter file; the cited inputs are kept as external pointers with their canonical Stacks tags.\par\medskip
\noindent\textit{Allowed standard prerequisites (editorial).}\begin{itemize}
\item Coverings by weakly contractible objects. Same chapter, the lemma that every object of the pro-etale site has a finite covering by weakly contractible objects, and the lemma producing coverings by affine weakly contractible objects for quasi-compact and quasi-separated objects. Supply the input coverings to which the hypercovering construction is applied.
\item Construction of hypercoverings from coverings with bounded index sets. Hypercoverings, the lemma producing a hypercovering with terms indexed by finite sets from coverings whose index sets are bounded by a cardinal. Supplies the hypercoverings K with finitely many weakly contractible terms in each degree.
\item Replacing a hypercovering by its degreewise disjoint unions. Hypercoverings, the remark that a hypercovering may be replaced by the hypercovering whose degree n term is the disjoint union of the terms in degree n. Supplies the passage to hypercoverings whose every term is a single object.
\item Continuity of the inclusion of the quasi-compact quasi-separated objects. Hypercoverings, the lemma that a continuous functor commuting with fibre products carries hypercoverings to hypercoverings. Supplies the passage from hypercoverings in the subcategory of quasi-compact and quasi-separated objects to hypercoverings in the pro-etale site.
\end{itemize}\medskip
\ifdefined\hypertarget\hypertarget{sourcepointer-1}{}\fi
\noindent\textit{Source reference.} Original lemma 0DAZ: sources/stacks-chow/tags.txt, source lines 16679--16679. Prerequisite pointer invoked at proof line(s) 3252; it names a standard fact of a sibling chapter that the counted argument uses, and every step of the counted proof built on it is retained verbatim in the delivered body.\par\medskip
\ifdefined\hypertarget\hypertarget{sourcepointer-2}{}\fi
\noindent\textit{Source reference.} Original lemma 094K: sources/stacks-chow/tags.txt, source lines 11553--11553. Prerequisite pointer invoked at proof line(s) 3234; it names a standard fact of a sibling chapter that the counted argument uses, and every step of the counted proof built on it is retained verbatim in the delivered body.\par\medskip
\ifdefined\hypertarget\hypertarget{sourcepointer-3}{}\fi
\noindent\textit{Source reference.} Original remark 0DB2: sources/stacks-chow/tags.txt, source lines 16682--16682. Prerequisite pointer invoked at proof line(s) 3240; it names a standard fact of a sibling chapter that the counted argument uses, and every step of the counted proof built on it is retained verbatim in the delivered body.\par\medskip
\ifdefined\hypertarget\hypertarget{sourcepointer-4}{}\fi
\noindent\textit{Source reference.} Original reference 098I: sources/stacks-proetale/pinned-original-source.tex, source lines 3087--3087. Prerequisite pointer invoked at proof line(s) 3255; it names a standard fact of the same chapter that the counted argument uses, and every step of the counted proof built on it is retained verbatim in the delivered body.\par\medskip
\ifdefined\hypertarget\hypertarget{sourcepointer-5}{}\fi
\noindent\textit{Source reference.} Original reference 0F4P: sources/stacks-proetale/pinned-original-source.tex, source lines 3143--3143. Prerequisite pointer invoked at proof line(s) 3233; it names a standard fact of the same chapter that the counted argument uses, and every step of the counted proof built on it is retained verbatim in the delivered body.\par\medskip
\begin{lemma}
\label{lemma-w-contractible-hypercovering}
Let $X$ be a scheme.
\begin{enumerate}
\item For every object $U$ of $X_\proetale$ there exists a hypercovering
$K$ of $U$ in $X_\proetale$ such that each term $K_n$ consists of a
single weakly contractible object of $X_\proetale$ covering $U$.
\item For every quasi-compact and quasi-separated object $U$ of $X_\proetale$
there exists a hypercovering $K$ of $U$ in $X_\proetale$ such that each
term $K_n$ consists of a single affine and weakly contractible object of
$X_\proetale$ covering $U$.
\end{enumerate}
\end{lemma}

\begin{proof}
Let $\mathcal{B} \subset \Ob(X_\proetale)$ be the set of weakly contractible
objects of $X_\proetale$. Every object $T$ of $X_\proetale$ has a
covering $\{T_i \to T\}_{i \in I}$ with $I$ finite and $T_i \in \mathcal{B}$ by
Lemma \ref{lemma-weakly-contractible-cover}.
By Hypercoverings, Lemma \ref{hypercovering-lemma-w-contractible}
we get a hypercovering $K$ of $U$ such that $K_n = \{U_{n, i}\}_{i \in I_n}$
with $I_n$ finite and $U_{n, i}$ weakly contractible.
Then we can replace $K$ by the hypercovering of $U$ given by $\{U_n\}$
in degree $n$ where $U_n = \coprod_{i \in I_n} U_{n, i}$
This is allowed by
Hypercoverings, Remark \ref{hypercovering-remark-take-unions-hypercovering-X}.

\medskip\noindent
Let $X_{qcqs, \proetale} \subset X_\proetale$ be the full subcategory
consisting of quasi-compact and quasi-separated objects.
A covering of $X_{qcqs, \proetale}$ will be a finite pro-\'etale covering.
Then $X_{qcqs, \proetale}$ is a site, has fibre products, and
the inclusion functor $X_{qcqs, \proetale} \to X_\proetale$ is continuous
and commutes with fibre products.
In particular, if $K$ is a hypercovering of an object $U$ in
$X_{qcqs, \proetale}$ then $K$ is a hypercovering of $U$ in $X_\proetale$
by Hypercoverings, Lemma
\ref{hypercovering-lemma-hypercovering-continuous-functor}.
Let $\mathcal{B} \subset \Ob(X_{qcqs, \proetale})$ be the set of
affine and weakly contractible objects. By 
Lemma \ref{lemma-get-many-weakly-contractible}
and the fact that finite unions of affines are affine,
for every object $U$ of $X_{qcqs, \proetale}$ there exists a covering
$\{V \to U\}$ of $X_{qcqs, \proetale}$ with $V \in \mathcal{B}$.
By Hypercoverings, Lemma \ref{hypercovering-lemma-w-contractible}
we get a hypercovering $K$ of $U$ such that $K_n = \{U_{n, i}\}_{i \in I_n}$
with $I_n$ finite and $U_{n, i}$ affine and weakly contractible.
Then we can replace $K$ by the hypercovering of $U$
given by $\{U_n\}$ in degree $n$ where
$U_n = \coprod_{i \in I_n} U_{n, i}$. This is allowed by
Hypercoverings, Remark \ref{hypercovering-remark-take-unions-hypercovering-X}.
\end{proof}

\end{document}
